\SourceLeaf{432}

also
\[
u=xyz
\]
Es ist folglich
\[
\begin{aligned}F(a+b)&=Fa+Fb\\[7pt]&\quad{}-\sin\operatorname{lemn}a.\sin\operatorname{lemn}b.\sin\operatorname{lemn}(a+b)\end{aligned}
\]
Setzt man
\[
Fa-\frac{a}{90^0}F90^0=Ga
\]
so ist
\[
\begin{aligned}G(a+b)&=Ga+Gb\\[7pt]&\quad{}-\sin\operatorname{lemn}a.\sin\operatorname{lemn}b.\sin\operatorname{lemn}(a+b)\end{aligned}
\]
\[
Ga=\frac{\mathrm dQa}{Qa.\mathrm da}
\]

\CurvaArticle{15}{[15.]}
\[
\begin{aligned}\frac{\varpi}{\pi}&=\Bigl\{1+\Bigl(\frac12\Bigr)^2.\frac12+
\Bigl(\frac{1.3}{2.4}\Bigr)^2.\frac14+
\\[7pt]&\quad{}\Bigl(\frac{1.3.5}{2.4.6}\Bigr)^2\frac18+..\Bigr\}\surd {\frac12}\end{aligned}
\]
\[
\begin{aligned}\frac{\varpi}{\pi}&=\Bigl(1+\frac{1.3}{4.4}.\frac19+
\frac{1.3.5.7}{4.4.8.8}.\frac1{81}+
\\[7pt]&\quad{}\frac{1.3.5.7.9.11}{4.4.8.8.12.12}.\frac1{729}+..\Bigr)\surd {\frac23}\end{aligned}
\]
\[
\begin{aligned}\frac{\varpi}{\pi}-\frac2{\varpi}&=\Bigl(\frac14.\frac13+\frac{1.3.5}{4.4.8}.\frac1{27}+
\\[7pt]&\quad{}\frac{1.3.5.7.9}{4.4.8.8.12}.\frac1{243}+
\frac{1.3.5.7.9.11.13}{4.4.8.8.12.12.16}.\frac1{2187}+..\Bigr)\surd {\frac23}\end{aligned}
\]

\CurvaArticle{16}{[16.]}
Ponendo
\[
\frac1{\mathrm M(1,\cos\varphi)}=N+a\cos2\varphi+b\cos4\varphi+\ .\ .\ .
\]
erit
\begingroup\small\[
\begin{array}{rl@{\hspace{12pt}}rl}
N&=1.393203&-\frac1\pi\log(1-\cos2\varphi)&=+0.220635\\
a&=0.581803&&+0.636620\cos2\varphi\\
b&=0.309601&&+0.318310\cos4\varphi\\
c&=0.209449&&+0.212207\cos6\varphi\\
d&=0.157960&&+0.159155\cos8\varphi\\
e&=0.126704.&&+0.127324\cos10\varphi
\end{array}
\]\endgroup

\SourceLeaf{433}

\[
N=\left(\frac{3.3.7.7.11.11.15}{2.4.6.8.10.12.14}\ .\ .\ .\right)^2
=2\frac{\varpi\varpi}{\pi\pi}
\]
\[
a=\frac12\left(\frac{5.5}{4.6}.\frac{9.9}{8.10}.\frac{13.13}{12.14}\ .\ .\ .\right)^2
=\frac4{\varpi\varpi}
\]
\[
b=\frac29N
\]
\[
\frac{\varpi}{\surd 8}=\frac{4.6.8.10.12.14}{5.5.9.9.13.13}\ .\ .\ .
\]
\[
\frac{\varpi}{2}=\frac32.\frac45.\frac76.\frac89.\frac{11}{10}.\frac{12}{13}\ .\ .\ .
\]

\CurvaArticle{17}{[17.]}
\[
\begin{aligned}
&1+\Bigl(\frac12\Bigr)^3x^4+\Bigl(\frac12.\frac34\Bigr)^3x^8+
\\[7pt]&\quad{}\Bigl(\frac12.\frac34.\frac56\Bigr)^3x^{12}+\ .\ .\ .\\[5pt]
&=\Bigl\{1+\Bigl(\frac14\Bigr)^2x^4+\Bigl(\frac14.\frac58\Bigr)^2x^8+
\\[7pt]&\quad{}\Bigl(\frac14.\frac58.\frac9{12}\Bigr)^2x^{12}+\ .\ .\Bigr\}^2
\end{aligned}
\]
\textit{Demonstratio.} Ponatur
\[
1+\left(\frac14\right)^2x^4+\left(\frac14.\frac58\right)^2x^8+\ .\ .\ .=t
\]
eritque, posito
\[
x(x^4-1)\frac{\mathrm{dd}t}{\mathrm dx^2}-(3x^4-1)\frac{\mathrm dt}{\mathrm dx}+x^3t=R,\quad R=0
\]
Hinc etiam
\[
\begin{aligned}&0=x\frac{\mathrm dR}{\mathrm dx}+R;\\[7pt]&\text{nec non}\quad
0=2xt\frac{\mathrm dR}{\mathrm dx}+2\left(t+3x\frac{\mathrm dt}{\mathrm dx}\right)R\end{aligned}
\]
unde fit, evolutione facta
\[
\begin{aligned}
0={}&xx(x^4-1)\left(2t\frac{\mathrm d^3t}{\mathrm dx^3}+6\frac{\mathrm dt}{\mathrm dx}\frac{\mathrm{dd}t}{\mathrm dx^2}\right)
\\[7pt]&\quad{}+3x(3x^4-1)\left(2t\frac{\mathrm{dd}t}{\mathrm dx^2}+2\frac{\mathrm dt}{\mathrm dx}\frac{\mathrm dt}{\mathrm dx}\right)\\[5pt]
&+(19x^4-1)2t\frac{\mathrm dt}{\mathrm dx}+8x^3tt
\end{aligned}
\]
sive ponendo $tt=u$
\[
\begin{aligned}0&=xx(x^4-1)\frac{\mathrm d^3u}{\mathrm dx^3}+3x(3x^4-1)\frac{\mathrm{dd}u}{\mathrm dx^2}
\\[7pt]&\quad{}+(19x^4-1)\frac{\mathrm du}{\mathrm dx}+8x^3u\end{aligned}
\]
cui aequationi invenitur respondere
\[
u=1+\left(\frac12\right)^3x^4+\left(\frac{1.3}{2.4}\right)^3x^8+\ .\ .
\]

\SourceLeaf{434}

Iam quum
\[
t=\frac{\surd 8}{\pi}\int\frac{\mathrm dz}{{}^{4}\!\surd{(1-z^4)(1-x^4z^4)}}
\]
$z=0$ usque ad $z=1$, fit, pro $x=1$,
\[
t=\frac{\surd 8}{\pi}.\frac{\varpi}{2}=\frac{\varpi}{\pi}\surd 2
\]
adeoque
\[
1+\left(\frac12\right)^3+\left(\frac{1.3}{2.4}\right)^3+
\left(\frac{1.3.5}{2.4.6}\right)^3+\ .\ .\ .=2\frac{\varpi\varpi}{\pi\pi}
\]
\begin{center}\textit{Idem alio modo}\end{center}
Valor seriei
\[
1+\left(\frac12\right)^3+\left(\frac{1.3}{2.4}\right)^3+
\left(\frac{1.3.5}{2.4.6}\right)^3+\ .\ .
\]
fit
\[
=\frac4{\pi\pi}\iint\frac{\mathrm d\varphi\,\mathrm d\psi}{\surd {(1-\cos\varphi^2\cos\psi^2)}}
\]
a $\varphi=0$ usque ad $\varphi=90^0$ et a $\psi=0$ usque ad $\psi=90^0$.
Faciendo itaque $\cos\varphi\cos\psi=\cos\upsilon$, idem valor fit
\[
=\frac4{\pi\pi}\iint\frac{\mathrm d\varphi\,\mathrm d\upsilon}{\surd {(\cos\varphi^2-\cos\upsilon^2)}}
\]
et quidem ab $\upsilon=0$ usque ad $\upsilon=90^0$ et a $\varphi=0$ usque ad $\varphi=\upsilon$.
Denique statuendo $\varphi+\upsilon=f$, $\upsilon-\varphi=g$, erit expressio nostra
\[
=\frac2{\pi\pi}\iint\frac{\mathrm df.\mathrm dg}{\surd {\sin f\sin g}}
\]
ab $g=0$ usque ad $g=90^0$ et a $f=g$ usque ad $f=180^0-g$.
Sed haud difficulter probatur integrale eundem valorem nancisci, si sumatur
ab $f=0$ usque ad $f=90^0$ et a $g=0$ usque ad $g=90^0$, unde ipsius valor deducitur
\[
=2\frac{\varpi\varpi}{\pi\pi}
\]
uti supra.

\SourceLeaf{435}

\par\medskip\Needspace{5\baselineskip}\pdfbookmark[0]{Sammlung von Rechnungen}{sammlung-rechnungen}\centerline{[17.]}

\begin{center}\textit{Sammlung von Rechnungen,}\end{center}
vornehmlich solchen, bei denen von meinen Methoden, die Factoren grosser
Zahlen zu finden, und von den \textsc{Wolframschen} Logarithmentafeln Gebrauch gemacht
ist.

\begin{center}\textit{Erste Rechnung für $e^{-\pi}=A$}\end{center}

Durch eine vorläufige Näherung war schon bekannt, dass $A$ bis auf die
$14^{\mathrm{te}}$ Figur $0,0432139182\ 6377$ sei, folglich $11.10^{10}A$ sehr genau $4753531008$
$=128.243.152827$. Es fragte sich also, ob die Zahl $152827$ sich noch in einfache
Factoren zerlegen lasse.

Die Division mit kleinen Primzahlen gelang nicht; man musste also zu
künstlichen Methoden seine Zuflucht nehmen. Hier fand sich nun, indem man
die Zahl selbst sowohl als verschiedene ihrer Vielfache mit den nächsten Quadraten
verglich, unter andern, dass
\[
552^2\equiv950,\quad677^2\equiv152\quad\mathrm{mod.}\ 152827
\]
woraus man sogleich schloss, dass
\[
1104^2\equiv3800,\quad3385^2\equiv3800
\]
und mithin die Zahl $152827$ keine Primzahl sei, weil sonst die Quadrate zweier
Zahlen, die beide kleiner als die Hälfte von jener sind, unmöglich congruent
sein könnten (\textit{Disqu. Arr. art.}\quad.). Durch die in diesem Werke gelehrte Methode
(\textit{art.}\quad) fanden sich nun die Factoren der Zahl $152827$, nemlich $67.2281$.

Man war also gewiss, dass sich der hyperbolische Logarithme von
\[
N=128.243.67.2281.11^{-1}.10^{-10}
\]
aus den vorhandenen Tafeln bestimmen lasse und zugleich dass derselbe von dem
gegebenen Logarithmen, $-\pi$, nur in der $10^{\mathrm{ten}}$ Decimalstelle abweichen könne.
Die zu dieser Differenz gehörige Absolutzahl brauchte also blos berechnet und
mit $N$ multiplicirt zu werden, um $A$ zu erhalten,
\[
A=Ne^{\log11+10\log10-\log2281-\log2144-\log972-\pi}=Ne^\delta
\]

\SourceLeaf{436}

Durch die \textsc{Wolframsche} Tafel war:

[Zur Abkürzung $\delta-\log(1+2135.10^{-13})=\delta'$,
$\delta'-\log(1+1443.10^{-17})=\delta''$ gesetzt.]

\begingroup\footnotesize\setlength{\arraycolsep}{1pt}\renewcommand{\arraystretch}{1.12}\[
\begin{array}{rcrlllll}
10\log10&&&&&&&\\[5pt]
&=&\GaussIIITranscriptionDigits{23,0258509299}&\GaussIIITranscriptionDigits{4045684017}&\GaussIIITranscriptionDigits{9914546843}&\GaussIIITranscriptionDigits{6420760110}&\GaussIIITranscriptionDigits{14886288}&\\[5pt]
\log11&=&\GaussIIITranscriptionDigits{2,3978952727}&\GaussIIITranscriptionDigits{9837054406}&\GaussIIITranscriptionDigits{1943577965}&\GaussIIITranscriptionDigits{1292998217}&\GaussIIITranscriptionDigits{06853937}&\\[5pt]
&&\GaussIIITranscriptionDigits{25,4237462027}&\GaussIIITranscriptionDigits{3882738424}&\GaussIIITranscriptionDigits{1858124808}&\GaussIIITranscriptionDigits{7713758327}&\GaussIIITranscriptionDigits{21740225}&\\[5pt]
l.2281&=&\GaussIIITranscriptionDigits{7,7323692222}&\GaussIIITranscriptionDigits{8438803081}&\GaussIIITranscriptionDigits{0466064812}&\GaussIIITranscriptionDigits{2168095619}&\GaussIIITranscriptionDigits{53159812}&\\[5pt]
l.2144&=&\GaussIIITranscriptionDigits{7,6704285221}&\GaussIIITranscriptionDigits{9069260675}&\GaussIIITranscriptionDigits{6232603654}&\GaussIIITranscriptionDigits{6055909444}&\GaussIIITranscriptionDigits{33575023}&\\[5pt]
l.972&=&\GaussIIITranscriptionDigits{6,8793558044}&\GaussIIITranscriptionDigits{6043907581}&\GaussIIITranscriptionDigits{0690427528}&\GaussIIITranscriptionDigits{9816593884}&\GaussIIITranscriptionDigits{53057834}&\\[5pt]
\pi&=&\GaussIIITranscriptionDigits{3,1415926535}&\GaussIIITranscriptionDigits{8979323846}&\GaussIIITranscriptionDigits{2643383279}&\GaussIIITranscriptionDigits{5028841971}&\GaussIIITranscriptionDigits{69399375}&\\[5pt]
&&\GaussIIITranscriptionDigits{25,4237462025}&\GaussIIITranscriptionDigits{2531295184}&\GaussIIITranscriptionDigits{0032479275}&\GaussIIITranscriptionDigits{3069440920}&\GaussIIITranscriptionDigits{09192044}&\\[5pt]
\delta&=&\GaussIIITranscriptionDigits{0,.........2}&\GaussIIITranscriptionDigits{1351443240}&\GaussIIITranscriptionDigits{1825645533}&\GaussIIITranscriptionDigits{4644317407}&\GaussIIITranscriptionDigits{12548181}&\\[5pt]
\delta'&=&\GaussIIITranscriptionDigits{0,..........}&\GaussIIITranscriptionDigits{...1443242}&\GaussIIITranscriptionDigits{4616770530}&\GaussIIITranscriptionDigits{2204949495}&\GaussIIITranscriptionDigits{65316893}&\\[5pt]
\delta''&=&\GaussIIITranscriptionDigits{0,..........}&\GaussIIITranscriptionDigits{.......242}&\GaussIIITranscriptionDigits{4616770634}&\GaussIIITranscriptionDigits{3329449495}&\GaussIIITranscriptionDigits{6431533124}&\\[5pt]
\frac12\delta''\delta''&=&\GaussIIITranscriptionDigits{0,..........}&\GaussIIITranscriptionDigits{..........}&\GaussIIITranscriptionDigits{..........}&\GaussIIITranscriptionDigits{.....29393}&\GaussIIITranscriptionDigits{8324222063}&\\[5pt]
e^{\delta''}&=&\GaussIIITranscriptionDigits{1,..........}&\GaussIIITranscriptionDigits{.......242}&\GaussIIITranscriptionDigits{4616770634}&\GaussIIITranscriptionDigits{3329478889}&\GaussIIITranscriptionDigits{4755755187}&\\[5pt]
e^{\delta'}&=&\GaussIIITranscriptionDigits{1,..........}&\GaussIIITranscriptionDigits{...1443242}&\GaussIIITranscriptionDigits{4616770634}&\GaussIIITranscriptionDigits{3679351089}&\GaussIIITranscriptionDigits{4781097632}&\\[5pt]
e^\delta&=&\GaussIIITranscriptionDigits{1,.........2}&\GaussIIITranscriptionDigits{1351443242}&\GaussIIITranscriptionDigits{4619851957}&\GaussIIITranscriptionDigits{0236156393}&\GaussIIITranscriptionDigits{8536512211}&\\[5pt]
11A&=&\GaussIIITranscriptionDigits{0,4753531009}&\GaussIIITranscriptionDigits{0149474751}&\GaussIIITranscriptionDigits{8595108889}&\GaussIIITranscriptionDigits{0081240330}&\GaussIIITranscriptionDigits{0920791693}&\GaussIIITranscriptionDigits{5}\\[5pt]
A&=&\GaussIIITranscriptionDigits{0,0432139182}&\GaussIIITranscriptionDigits{6377224977}&\GaussIIITranscriptionDigits{4417737171}&\GaussIIITranscriptionDigits{7280112757}&\GaussIIITranscriptionDigits{2810981063}&
\end{array}
\]\endgroup

Um sich von der Richtigkeit dieses Resultats durch eine zweite Rechnung
zu versichern, multiplicirte man die Zahl $A$ durch $599.10^8$, wodurch sich ergab
$2588513703,999957..$, so dass man also eine sehr leichte Rechnung übrig hatte,
wenn die Zahl $2588513704$ sich in Factoren kleiner als $10000$ zerlegen liess.
Nach angestelltem Versuch fand sich $2588513704=8.7.17.2719027$. Es kam
also darauf an, ob $2719027$ eine Primzahl sei. Man fand, dass $-1848$ gewiss
ein quadratischer Rest von $2719027$ sein müsse, wenn diese Zahl eine Primzahl
sei, und dass sie in diesem Fall \textit{einmal} unter der Form $3xx+616yy$ enthalten
sein müsse und umgekehrt, dass sie durch diese Form entweder gar nicht oder
\textit{mehr als einmal} müsse dargestellt werden können, wenn sie zusammengesetzt sei.
Allein die Exclusionsmethode lehrte, dass jene Zahl wirklich nur einmal unter
der Form $3xx+616yy$ enthalten sei, nemlich $2719027=3.197^2+616.65^2$,
\SourceLeaf{437}
woraus also mit Gewissheit folgte, dass $2719027$ eine Primzahl und folglich die
versuchte Methode diesmal nicht anwendbar sei.

Nach einigen andern vergeblichen Versuchen kam man endlich auf folgenden,
der besser gelang. Die Zahl $A$ multiplicirt mit $10^{14}$ gab $4321391826377,25$
und die Zahl $4321391826375$ zerfiel in die Factoren $125.81.13.32831087$. Dass
die Zahl $32831087$ keine Primzahl sei, folgte daraus, dass sie nicht unter der
Form $xx+190yy$ enthalten war; man wandte also die erste der in den \textit{Disqui. Arr.}
gelehrten Methoden darauf an, wodurch sich entdeckte, dass sie das Product
aus $373.8819$ sei.

\begin{center}\textit{Zweite Rechnung für $e^{-\pi}=A$}\end{center}

$59.10^9A$ fand sich sehr genau $=2549621178=54.13.3631939$. Die Zahl
$3631939$ zerfiel in die Factoren $1091.3329$ (welche aus
\[
3631939=40.232^2+19.279^2=40.300^2+19.41^2
\]
gefunden waren). Also
\[
\begin{aligned}59A&=702.1091.3329.10^{-9}\\[7pt]&\quad{}e^{9\log10+\log59-\log702-\log1091-\log3329-\pi}=Ne^\delta\end{aligned}
\]
Aus \textsc{Wolframs} Tafel fand sich [$\delta-\log(1-17157.10^{-14})=\delta'$,
$\frac12\delta'\delta'+\frac16\delta'^3=\varepsilon$ gesetzt]
\begingroup\footnotesize\setlength{\arraycolsep}{1pt}\renewcommand{\arraystretch}{1.12}\[
\begin{array}{rcrlllll}
C9l.10&=&\GaussIIITranscriptionDigits{21,2767341630}&\GaussIIITranscriptionDigits{5358884383}&\GaussIIITranscriptionDigits{8076907840}&\GaussIIITranscriptionDigits{7221315900}&\GaussIIITranscriptionDigits{86602341}&\\[5pt]
C.l.59&=&\GaussIIITranscriptionDigits{5,9224625560}&\GaussIIITranscriptionDigits{9428054938}&\GaussIIITranscriptionDigits{3949626280}&\GaussIIITranscriptionDigits{3023759366}&\GaussIIITranscriptionDigits{53210670}&\\[5pt]
l.702&=&\GaussIIITranscriptionDigits{6,5539334040}&\GaussIIITranscriptionDigits{2581111965}&\GaussIIITranscriptionDigits{6455273791}&\GaussIIITranscriptionDigits{0722868232}&\GaussIIITranscriptionDigits{39752589}&\\[5pt]
l.1091&=&\GaussIIITranscriptionDigits{6,9948499858}&\GaussIIITranscriptionDigits{3307081851}&\GaussIIITranscriptionDigits{1895817110}&\GaussIIITranscriptionDigits{3840806982}&\GaussIIITranscriptionDigits{08318537}&\\[5pt]
l.3329&=&\GaussIIITranscriptionDigits{8,1104272375}&\GaussIIITranscriptionDigits{7502494295}&\GaussIIITranscriptionDigits{2021653586}&\GaussIIITranscriptionDigits{1576658324}&\GaussIIITranscriptionDigits{67001597}&\\[5pt]
\pi&=&\GaussIIITranscriptionDigits{3,1415926535}&\GaussIIITranscriptionDigits{8979323846}&\GaussIIITranscriptionDigits{2643383279}&\GaussIIITranscriptionDigits{5028841971}&\GaussIIITranscriptionDigits{69399375}&\\[5pt]
-\delta&=&\GaussIIITranscriptionDigits{0,.........1}&\GaussIIITranscriptionDigits{7156951280}&\GaussIIITranscriptionDigits{5042661888}&\GaussIIITranscriptionDigits{1414250778}&\GaussIIITranscriptionDigits{24285109}&\\[5pt]
\delta'&=&\GaussIIITranscriptionDigits{0,..........}&\GaussIIITranscriptionDigits{.....48720}&\GaussIIITranscriptionDigits{9675470563}&\GaussIIITranscriptionDigits{5420349120}&\GaussIIITranscriptionDigits{2333831382}&\\[5pt]
\varepsilon&=&\GaussIIITranscriptionDigits{0,..........}&\GaussIIITranscriptionDigits{..........}&\GaussIIITranscriptionDigits{..........}&\GaussIIITranscriptionDigits{1186866339}&\GaussIIITranscriptionDigits{3606594227}&\\[5pt]
e^{\delta'}&=&\GaussIIITranscriptionDigits{1,..........}&\GaussIIITranscriptionDigits{.....48720}&\GaussIIITranscriptionDigits{9675470563}&\GaussIIITranscriptionDigits{6607215459}&\GaussIIITranscriptionDigits{5940425609}&\\[5pt]
1-e^\delta&=&\GaussIIITranscriptionDigits{0,.........1}&\GaussIIITranscriptionDigits{7156951279}&\GaussIIITranscriptionDigits{0324613026}&\GaussIIITranscriptionDigits{9032989386}&\GaussIIITranscriptionDigits{4786573955}&\\[5pt]
59A&=&\GaussIIITranscriptionDigits{2,5496211775}&\GaussIIITranscriptionDigits{6256273669}&\GaussIIITranscriptionDigits{0646493131}&\GaussIIITranscriptionDigits{9526652679}&\GaussIIITranscriptionDigits{5847882752}&\GaussIIITranscriptionDigits{6}\\[5pt]
A&=&\GaussIIITranscriptionDigits{0,0432139182}&\GaussIIITranscriptionDigits{6377224977}&\GaussIIITranscriptionDigits{4417737171}&\GaussIIITranscriptionDigits{7280112757}&\GaussIIITranscriptionDigits{2810981063}&\GaussIIITranscriptionDigits{6}
\end{array}
\]\endgroup

\SourceLeaf{438}

Zugleich kann man aus der Übereinstimmung beider Rechnungen schliessen,
dass die Logarithmen von $2,3,5,11,13,59,67,1091,2281,3329$ in der
\textsc{Wolframschen} Tafel bis auf die letzte Figur richtig sind.

\begin{center}\textit{Doppelte Berechnung von $e^{-\frac14\pi}=A$.}\end{center}
\begin{center}Erste Rechnung.\end{center}

Da man durch eine schon vorher angestellte Rechnung wusste, dass der
Werth von $A$ bis auf die letzte Zifer $=0,4559381277\ 6599$ sei, so war
$19.131.A$ oder $2489A=1134,8300000095$ oder $248900A$ sehr genau $=113483$.
Man fand $113483=283.401$ aus
\[
113483=311^2+58.17^2=79^2+58.43^2
\]
und hiemit
\[
\begin{aligned}2489A&=113483.10^{-2}.\\[7pt]&\quad{}e^{2\log10-\log283-\log401+\log2489-\frac14\pi}=Ne^\delta\end{aligned}
\]
Aus \textsc{Wolframs} Tafeln erhielt man

[wenn zur Abkürzung
\[
\delta-\log(1+8428.10^{-15})=\delta'
\]
gesetzt wird]:
\begingroup\footnotesize\setlength{\arraycolsep}{1pt}\renewcommand{\arraystretch}{1.12}\[
\begin{array}{rcrlllll}
\delta&=&\GaussIIITranscriptionDigits{0,..........}&\GaussIIITranscriptionDigits{.842825208}&\GaussIIITranscriptionDigits{2107272461}&\GaussIIITranscriptionDigits{7421888198}&\GaussIIITranscriptionDigits{29069513}&\\[5pt]
\delta'&=&\GaussIIITranscriptionDigits{0,..........}&\GaussIIITranscriptionDigits{.....25208}&\GaussIIITranscriptionDigits{2142788053}&\GaussIIITranscriptionDigits{7419892695}&\GaussIIITranscriptionDigits{5615344103}&\\[5pt]
\frac12\delta'\delta'&=&&&&\GaussIIITranscriptionDigits{\phantom{0}317727033}&\GaussIIITranscriptionDigits{5630835760}&\\[5pt]
\frac12\delta'^3&=&&&&&\GaussIIITranscriptionDigits{\phantom{0000000}267}&\\[5pt]
e^{\delta'}&=&\GaussIIITranscriptionDigits{1,..........}&\GaussIIITranscriptionDigits{.....25208}&\GaussIIITranscriptionDigits{2142788053}&\GaussIIITranscriptionDigits{7737619729}&\GaussIIITranscriptionDigits{1246180130}&\\[5pt]
e^\delta&=&\GaussIIITranscriptionDigits{1,..........}&\GaussIIITranscriptionDigits{.842825208}&\GaussIIITranscriptionDigits{2142790178}&\GaussIIITranscriptionDigits{3220163906}&\GaussIIITranscriptionDigits{2965706721}&\\[5pt]
2489A&&&&&&&\\[5pt]
&=&\GaussIIITranscriptionDigits{1134,8300000095}&\GaussIIITranscriptionDigits{6463331037}&\GaussIIITranscriptionDigits{8102578065}&\GaussIIITranscriptionDigits{2249279282}&\GaussIIITranscriptionDigits{5372958193}&\\[5pt]
A&=&\GaussIIITranscriptionDigits{0,4559381277}&\GaussIIITranscriptionDigits{6599623676}&\GaussIIITranscriptionDigits{5921294728}&\GaussIIITranscriptionDigits{0294194166}&\GaussIIITranscriptionDigits{04366}&
\end{array}
\]\endgroup

\SourceLeaf{439}

\begin{center}\emph{Zweite Rechnung für $e^{-\frac14\pi}=A$}\end{center}

$113A$ wurde gefunden $51{,}52100843755$. Die Zahl $515210084352$ zerfiel in die Factoren $27.131072.145583$.
Endlich erhielt man aus
\[
4.145583=437^2+163.49^2=763^2+163
\]
$145583=197.739$ mithin
\[
\begin{aligned}113A&=96^3.788.739.10^{-10}.\\[7pt]&\quad{}e^{10\log10+\log113-3\log96-\log788-\log739-\frac14\pi}=Ne^\delta\end{aligned}
\]
[und wenn $\delta-\log(1+4576.10^{-15})=\delta'$ gesetzt wird]

\begingroup\footnotesize\setlength{\arraycolsep}{1pt}\renewcommand{\arraystretch}{1.12}\[
\begin{array}{r@{\,=\,}r@{\;}l@{\hspace{2pt}}l@{\hspace{2pt}}l@{\hspace{2pt}}l@{\hspace{2pt}}l@{\hspace{2pt}}l}
\delta & 0{,} & .\,.\,.\,.\,.\,.\,.\,.\,.\,. & 4575948387 & 0747214839 & 8399455217 & 14190720 & \\[5pt]
-\delta' & 0{,} & .\,.\,.\,.\,.\,.\,.\,.\,.\,. & .\,.\,.\,.\,.\,51612 & 8205796360 & 1919946163 & 4337976121 & \\[5pt]
+\tfrac12\delta'\delta' & 0{,} &  &  &  & 1331941624 & 0928492341 & \\[5pt]
1-e^{\delta'} & 0{,} & .\,.\,.\,.\,.\,.\,.\,.\,.\,. & .\,.\,.\,.\,.\,51612 & 8205796360 & 0588043539 & 3409483780 & \\[5pt]
e^\delta & 1{,} & .\,.\,.\,.\,.\,.\,.\,.\,.\,. & 4575978387 & 1794180021 & 9145023046 & 2961445343 & \\[5pt]
113A & 51{,} & 5210084375 & 5757475454 & 9106304267 & 3243940900 & 3220627240 & 5\\[5pt]
A & 0{,} & 4559381277 & 6599623676 & 5921294728 & 0294194166 & 0436523820 & \\[5pt]
\end{array}
\]\endgroup

\[
e^{-\frac14\pi}=0{,}4559382=3671.54.23.10^{-7}
\]

\begin{center}\emph{Dritte Rechnung für $e^{-\frac14\pi}=A$.}\end{center}

$455938128=144.3166237$ Die Divisoren der Zahl $3166237$ fand man, indem man die Periode der reducirten Form
$(1,1779,-1396)$ entwickelte, in welcher die Form $(-1897,530,1521)$ die sechste war. Hieraus
\[
\begin{aligned}(11+28.1779)^2\equiv1521&=39^2,\\[6pt]&\text{und }3166237=107.127.233\end{aligned}
\]
\[
\begin{aligned}A&=455938128.10^{-9}.\\[7pt]&\quad{}e^{9\log10-\frac14\pi-\log455938128}=Ne^\delta\end{aligned}
\]
[und $\delta-\log(1-513.10^{-12})=\delta'$ gesetzt]:

\begingroup\footnotesize\setlength{\arraycolsep}{1pt}\renewcommand{\arraystretch}{1.12}\[
\begin{array}{r@{\,=\,}r@{\;}l@{\hspace{2pt}}l@{\hspace{2pt}}l@{\hspace{2pt}}l@{\hspace{2pt}}l@{\hspace{2pt}}l}
-\delta & 0{,} & .\,.\,.\,.\,.\,.\,.\,.\,.\,5 & 1323578556 & 7221212746 & 8124092623 & 9185579000 & \\[5pt]
-\delta' & 0{,} & .\,.\,.\,.\,.\,.\,.\,.\,.\,. & .\,.\,23578543 & 5636712701 & 8105102450 & 7737514264 & \\[5pt]
+\tfrac12\delta'\delta' &  &  &  & \phantom{00000}27797 & 3858291971 & 9406911797 & \\[5pt]
-\tfrac16\delta^{\prime3} &  &  &  &  & \phantom{00000000}21 & 8473957577 & \\[5pt]
1-e^{\delta'} & 0{,} & .\,.\,.\,.\,.\,.\,.\,.\,.\,. & .\,.\,23578543 & 5636684904 & 4246810500 & 6804560044 & \\[5pt]
1-e^\delta & 0{,} & .\,.\,.\,.\,.\,.\,.\,.\,.\,5 & 1323578543 & 5515726975 & 9436016940 & 9818422176 & \\[5pt]
\end{array}
\]\endgroup

\par\Needspace{8\baselineskip}\SourceLeaf{440}

\begin{center}\emph{Berechnung von $e^{-\frac94\pi}=A$.}\end{center}

Durch Näherung war bereits gefunden $A=0{,}0008514383\ 42805$. Nun fand sich
$8514383436=4.9.13.19.307.3119$, mithin sehr genau
\[
A=4.9.13.19.307.3119.10^{-13}
\]
aus \textsc{Wolframs} Tafel
\[
13\log10-\log3119-\log1842-\log1482-\frac94\pi=\delta
\]
[und
\[
\begin{aligned}\delta-\log(1-9335.10^{-13})&=\delta',\\[7pt]&
\delta'-\log(1-285.10^{-16})=\delta''\end{aligned}
\]
gesetzt]:

\begingroup\footnotesize\setlength{\arraycolsep}{1pt}\renewcommand{\arraystretch}{1.12}\[
\begin{array}{r@{\,=\,}r@{\;}l@{\hspace{2pt}}l@{\hspace{2pt}}l@{\hspace{2pt}}l@{\hspace{2pt}}l@{\hspace{2pt}}l}
-\delta & 0{,} & .\,.\,.\,.\,.\,.\,.\,.\,.\,9 & 3352850560 & 8342583868 & 5326823995 & 3184946898 & 8\\[5pt]
-\delta' & 0{,} & .\,.\,.\,.\,.\,.\,.\,.\,.\,. & .\,.\,.\,2850517 & 2631458732 & 9539039720 & 5781872612 & 0\\[5pt]
-\delta'' & 0{,} &  & .\,.\,.\,.\,.\,.\,.\,517 & 2631458326 & 8289039720 & 5396053412 & 0\\[5pt]
\tfrac12\delta''\delta'' & 0{,} &  &  &  & \phantom{0000}133780 & 5810183616 & 8\\[5pt]
1-e^{\delta''} & 0{,} &  & .\,.\,.\,.\,.\,.\,.\,517 & 2631458326 & 8288905939 & 9585869795 & 2\\[5pt]
1-e^{\delta'} & 0{,} &  & .\,.\,.\,2850517 & 2631458326 & 6814705974 & 3354407457 & 0\\[5pt]
1-e^\delta & 0{,} & .\,.\,.\,.\,.\,.\,.\,.\,.\,9 & 3352850517 & 2604848748 & 0300042494 & 7639127186 & 6\\[5pt]
A & 0{,} & 0008514383 & 4280515803 & 5852453295 & 4846487994 & 1872486024 & \\[5pt]
\multicolumn{2}{c}{}&\multicolumn{6}{r}{8176915}\\[5pt]
\end{array}
\]\endgroup

\[
e^{-\frac94\pi}\quad\text{satis exacte}\quad=4.11.13.29.3593.7^{-1}.10^{-10}
\]

\begin{center}\emph{Ecce iam computatum pro $e^{\frac12\pi}$.}\end{center}

Per logarithmos brigg. invenimus praeter propter $e^{\frac12\pi}=4{,}810484$.
Est vero $48104847=2293.37.7.81$ et
\[
\begin{gathered}
-\log81.259.2293+7\log10+\tfrac12\pi\\[5pt]=-0{,}.\,.\,.\,.\,.\,15214\ 7666454820\ 0537824776\ 3190\\[5pt]
\mathrm{num}\log=1-0{,}.\,.\,.\,.\,.\,15214\ 7550690316\ 7468363738\ 6798\\[5pt]
e^{\frac12\pi}=4{,}8104773809\ 6535165547\ 3044648993\ 1536
\end{gathered}
\]

\par\Needspace{8\baselineskip}\SourceLeaf{441}

\begin{center}\emph{Computus secundus pro $e^{\frac12\pi}$.}\end{center}

Invenimus $48104773808=195497.13^3.7.16$ superest itaque ut utrum sit $195497$ numerus non primus investigetur.
Tentamus itaque aequationem
\[
195497=16xx+\square
\]
Lim $x=111$, Excl. valores $x=\text{---\ ---\ ---}$
\[
\begin{aligned}
195497&=44^2+193561=76^2+189721\\[6pt]&=356^2+68761=364^2+63001\\[5pt]
&=364^2+251^2\quad\text{quare }195497\text{ est prim.}
\end{aligned}
\]
\[
\frac{26905}{5593}=3+\tfrac47+\tfrac{16}{17}+\tfrac{14}{47}=4{,}8104773824\ 4\ .\ .
\]
\[
\begin{gathered}
\log\frac{7.799}{5.5381}+\tfrac12\pi\\[5pt]=-0{,}.\,.\,.\,.\,.\,.\,.\,.\,.\,3\ 0703542806\ 9410475204\ 9155\\[5pt]
\mathrm{num}\log=1-0{,}.\,.\,.\,.\,.\,.\,.\,.\,.\,3\ 0703542802\ 2275098164\ 7660
\end{gathered}
\]

\begin{center}\emph{Computus pro $e^{-\frac12\pi}$}\end{center}

\[
\begin{aligned}e^{-\frac12\pi}\quad\text{praeter propter}&=0{,}20787957\\[6pt]&=7151.19.17.9.10^{-8}\end{aligned}
\]
\[
\begin{gathered}
-\log51.57.7151+8\log10-\tfrac12\pi\\[5pt]=-0{,}.\,.\,.\,.\,.\,.\,.\,305\ 5019697961\ 8740193606\ 8277\\[5pt]
\mathrm{num}\log=1-0{,}.\,.\,.\,.\,.\,.\,.\,305\ 5019651296\ 1477199011\ 4138
\end{gathered}
\]
Investigatio divisorum numeri $2078795763=99.20997937{,}005$
\begin{center}---\quad---\quad---\end{center}
\[
20997937=1848.34^2+4343^2\quad\text{numerus primus}
\]

\begin{center}\emph{Computus secundus pro $e^{-\frac12\pi}$}\end{center}

\[
11.e^{-\frac12\pi}=2{,}286675339858378
\]
\[
228667536=81.16.176441
\]
\[
176441=880xx+yy=73.2417
\]
\[
\begin{gathered}
\log\frac{11}{144.657.2417}+8\log10-\tfrac12\pi\\[5pt]=-0{,}.\,.\,.\,.\,.\,.\,.\,.\,88\ 0825474705\ 3269478285\ 7110
\end{gathered}
\]
\par\nobreak\medskip\centerline{\rule{.44\linewidth}{.3pt}}
